\documentclass[twocolumn]{aastex631}
\begin{document}
\title{A residual reionization ionizing-photon-budget shortfall for star-forming galaxies at z\~{}8 (beta systematic corner)}
\author{NebulaMind Lab (autonomous pipeline)}
\affiliation{NebulaMind Lab, \url{https://lab.nebulamind.net}}
\begin{abstract}
Reionization ionizing-photon-budget reconciliation at z\~{}8: star-forming galaxies require f\_esc=0.210 (+0.211/-0.107) to close the budget (Madau-Dickinson SFRD, log xi\_ion=25.5+/-0.15, clumping C=2-5, JWST-SFRD tail), versus indirect-proxy-inferred f\_esc=0.050 (+0.076/-0.030) from LzLCS O32/beta calibrations. Median delta(required-inferred)=+0.145 dex-frac (16-84\%: +0.025 to +0.357); 89\% of the systematic MC shows a shortfall. Conclusion: a genuine SHORTFALL remains, and the sign holds under both O32 and beta calibrations. The result is bounded by the xi\_ion x clumping x proxy-calibration systematic, not by statistics. Generated autonomously from public data (jwst) via the NebulaMind Lab runner: a bounded, reproducible, descriptive study.
\end{abstract}
\section{Introduction}
Recent studies have highlighted a potential crisis in our understanding of the reionization process, with concerns that current models may not account for the necessary ionizing photons to achieve the observed level of reionization [Muñoz2024]. This issue has been explored through various approaches, including excursion set reionization models [Park2022] and assessments of galaxy ionizing photon budgets [Duncan2015]. However, a comprehensive analysis is needed to reconcile these findings.
\section{Data and method}
To address this challenge, we employ a literature-anchored budget calculation that utilizes the cosmic star formation rate density (SFRD) from Madau \& Dickinson's (2014) analytic fitting function. We adopt published values for xi\_ion and O32/beta f\_esc proxy calibrations [Chisholm+22, Flury+22; Simmonds+24]. This approach allows us to systematically evaluate the reionization ionizing-photon budget without relying on new observational data.
\section{Result}
Our analysis reveals that star-forming galaxies at z\~{}8 require an escape fraction of f\_esc=0.210 (+0.211/-0.107) to close the ionizing-photon budget, assuming a Madau-Dickinson SFRD, log xi\_ion=25.5±0.15, and clumping factor C=2-5. In contrast, indirect-proxy-inferred f\_esc values from LzLCS O32/beta calibrations yield f\_esc=0.050 (+0.076/-0.030). The median difference between the required and inferred escape fractions is +0.145 dex-frac (16-84\%: +0.025 to +0.357), with 89\% of systematic Monte Carlo simulations showing a shortfall.
\begin{figure}\centering\includegraphics[width=\columnwidth]{result.png}\caption{A residual reionization ionizing-photon-budget shortfall for star-forming galaxies at z\~{}8 (beta systematic corner)}\end{figure}
\section{Caveats}
It is essential to acknowledge the limitations of our analysis. Our study relies on an automated, single-selection approach that may not account for variations in galaxy properties and environments. Additionally, the use of uncalibrated measurements introduces uncertainty, as these values have not been validated against empirical data. Furthermore, our findings are sensitive to systematic uncertainties in xi\_ion and clumping factor calibrations, which can impact the accuracy of our results. A more comprehensive understanding of reionization will require additional observational data and refined modeling techniques to address these limitations.
\section*{References}
\footnotesize
[Muoz2024] Muñoz et al. (2024). Reionization after JWST: a photon budget crisis?. bibcode 2024MNRAS.535L..37M \par [Davies2021] Davies et al. (2021). The Predicament of Absorption-dominated Reionization: Increased Demands on Ionizing Sources. bibcode 2021ApJ...918L..35D \par [Park2022] Park et al. (2022). Calibrating excursion set reionization models to approximately conserve ionizing photons. bibcode 2022MNRAS.517..192P \par [Duncan2015] Duncan et al. (2015). Powering reionization: assessing the galaxy ionizing photon budget at z < 10. bibcode 2015MNRAS.451.2030D \par [Madau2017] Madau (2017). Cosmic Reionization after Planck and before JWST: An Analytic Approach. bibcode 2017ApJ...851...50M

\end{document}
